% !TEX TS-program = pdflatex
% !TEX encoding = UTF-8 Unicode

% This is a simple template for a LaTeX document using the "article" class.
% See "book", "report", "letter" for other types of document.

\documentclass[12pt]{article} % use larger type; default would be 10pt
\usepackage[fontsize=12pt]{fontsize}
\renewcommand{\baselinestretch}{1.5} % linespacing=1.5
\usepackage[T2A]{fontenc} % кодировка
\usepackage{fontspec} % for changing fonts 
% \setmainfont{Times New Roman}
\setmainfont{Minion Pro}
% \setmainfont{Roboto}
% \usepackage[utf8]{inputenc} % set input encoding (not needed with XeLaTeX)
\usepackage[english]{babel} % for cyrillic letters support
\usepackage[parfill]{parskip} % remove space vefore parapraph
\setlength{\parskip}{0pt} % additional skip between paragraphs


%%% PAGE DIMENSIONS
\usepackage{geometry} % to change the page dimensions
\geometry{a4paper} % or letterpaper (US) or a5paper or....
\geometry{left=30mm,right=20mm,top=15mm, bottom=15mm} % for example, change the margins to 2 inches all round
% \graphicspath{{./images/}}
% \geometry{landscape} % set up the page for landscape
%   read geometry.pdf for detailed page layout information

\usepackage{graphicx} % support the \includegraphics command and options

% \usepackage[parfill]{parskip} % Activate to begin paragraphs with an empty line rather than an indent

%%% PACKAGES
\usepackage{blindtext}
\usepackage{booktabs} % for much better looking tables
\usepackage{multirow}
\usepackage{array} % for better arrays (eg matrices) in maths
\usepackage{paralist} % very flexible & customisable lists (eg. enumerate/itemize, etc.)
\usepackage{verbatim} % adds environment for commenting out blocks of text & for better verbatim
\usepackage{subfig} % make it possible to include more than one captioned figure/table in a single float
\usepackage{cmap} % search in pdf
\usepackage{longtable} % for long tables 
\usepackage{lscape}
\usepackage{hyperref} % for hyperlinks
\usepackage{listings} % for code listings
% These packages are all incorporated in the memoir class to one degree or another...

%% Useful packages
\usepackage[colorinlistoftodos]{todonotes}
% \usepackage[colorlinks=true, allcolors=blue]{hyperref}



\usepackage{amsmath,amsthm,amsfonts,amssymb,amscd, fancyhdr, color, comment, graphicx, environ}
\usepackage{float}
\usepackage{mathrsfs}
\usepackage[math-style=ISO]{unicode-math}
\setmathfont{TeX Gyre Termes Math}
\setmonofont{Hack Nerd Font Mono}
\usepackage{lastpage}
\usepackage[dvipsnames]{xcolor}
% \usepackage[framemethod=TikZ]{mdframed}
\usepackage{indentfirst}
\usepackage{thmtools}
\usepackage{shadethm}
\usepackage{setspace}


\definecolor{codegreen}{rgb}{0,0.6,0}
\definecolor{codegray}{rgb}{0.5,0.5,0.5}
\definecolor{codepurple}{rgb}{0.58,0,0.82}
\definecolor{backcolour}{rgb}{0.95,0.95,0.92}


%%% HEADERS & FOOTERS
\usepackage{fancyhdr} % This should be set AFTER setting up the page geometry
\pagestyle{fancy} % options: empty , plain , fancy
\renewcommand{\headrulewidth}{0pt} % customise the layout...
\lhead{}\chead{}\rhead{}
\lfoot{}\cfoot{\thepage}\rfoot{}

%%% SECTION TITLE APPEARANCE
\usepackage{titlesec}
\usepackage{sectsty}
\sectionfont{\centering}
\subsectionfont{\centering}

\usepackage{enumitem}
\setlist{nolistsep}

\hypersetup{
    colorlinks=true,
    linkcolor=black,
    filecolor=magenta,      
    urlcolor=blue,
    pdftitle={Overleaf Example},
    pdfpagemode=FullScreen,
    }
\urlstyle{same}




\lstdefinestyle{mystyle}{
    extendedchars=\true,
    inputencoding=utf8x,
    backgroundcolor=\color{backcolour},   
    commentstyle=\color{codegreen},
    keywordstyle=\color{magenta},
    numberstyle=\tiny\color{codegray},
    stringstyle=\color{codepurple},
    basicstyle=\ttfamily\footnotesize,
    breakatwhitespace=false,         
    breaklines=true,                 
    captionpos=b,                    
    keepspaces=true,                 
    numbers=left,                    
    numbersep=5pt,                  
    showspaces=false,                
    showstringspaces=false,
    showtabs=false,                  
    tabsize=2
}

\lstset{style=mystyle}




% \allsectionsfont{\sffamily\mdseries\upshape} % (See the fntguide.pdf for font help)
% (This matches ConTeXt defaults)

%%% ToC (table of contents) APPEARANCE
% \usepackage[nottoc,notlof,notlot]{tocbibind} % Put the bibliography in the ToC
% \usepackage[titles,subfigure]{tocloft} % Alter the style of the Table of Contents
% \renewcommand{\cftsecfont}{\rmfamily\mdseries\upshape}
% \renewcommand{\cftsecpagefont}{\rmfamily\mdseries\upshape} % No bold!
\makeatletter
\renewcommand{\fnum@figure}{Picture. \thefigure}
\makeatother



\renewcommand\lstlistingname{Algorithm}
\renewcommand\lstlistlistingname{Algorithms}
\def\lstlistingautorefname{Alg.}


\title{Тестрование Авторизции и регистрации}
\author{Бактыбеков Н.Б.}

\begin{document}
\begin{titlepage}

\newcommand{\HRule}{\rule{\linewidth}{0.5mm}} % Defines a new command for the horizontal lines, change thickness here

%----------------------------------------------------------------------------------------
%	LOGO SECTION
%----------------------------------------------------------------------------------------
\centering
\includegraphics[width=8cm]{title/logo_inai.jpg}\\[1cm] % Include a department/university logo - this will require the graphicx package
 
%----------------------------------------------------------------------------------------

\center % Center everything on the page

%----------------------------------------------------------------------------------------
%	HEADING SECTIONS
%----------------------------------------------------------------------------------------

\textsc{\LARGE Проверка авторизации и регистрации}\\[1.5cm] 
\textsc{\Large }\\[0.5cm] 
\textsc{\large WIN-1-22}\\[0.5cm] 

%----------------------------------------------------------------------------------------
%	TITLE SECTION
%----------------------------------------------------------------------------------------
\makeatletter
\HRule \\[0.4cm]
{ \huge \bfseries \@title}\\[0.4cm] % Title of your document
\HRule \\[6.5cm]
 
%----------------------------------------------------------------------------------------
%	AUTHOR SECTION
%----------------------------------------------------------------------------------------

\begin{minipage}{0.4\textwidth}
\begin{flushleft} \large
\emph{Выполнил:}\\
\@author % Your name



\end{flushleft}
\end{minipage}
~
\begin{minipage}{0.4\textwidth}
\begin{flushright} \large
\emph{} \\
.\\[1.2em] 
\end{flushright}
\end{minipage}\\[2cm]
\makeatother

%----------------------------------------------------------------------------------------
%	DATE SECTION
%----------------------------------------------------------------------------------------

{\large \today}\\[2cm] % Date, change the \today to a set date if you want to be precise

\vfill % Fill the rest of the page with whitespace

\end{titlepage}


% \maketitle

\newpage
% \tableofcontents{}
\setcounter{page}{1}

% \newpage
\addcontentsline{toc}{section}{Проверка авторизации и регистрации}


\section*{Проверка авторизации и регистрации}


\textbf{Введение}

Авторизационное и регистрационное тестирование — это тип тестирования программного обеспечения, который гарантирует, что пользователи могут получить доступ только к тем ресурсам и функциям, на доступ к которым им предоставлен доступ. Этот тип тестирования важен для обеспечения безопасности и конфиденциальности данных пользователей.

\textbf{Цели тестирования}

Целями авторизационного и регистрационного тестирования являются:

\begin{compactitem}
    \item    Уьбедиться, что пользователи могут получить доступ только к тем ресурсам и функциям, к которым у них есть доступ.
    \item    Уьбедиться, что пользователи не могут получить доступ к ресурсам и функциям, на доступ к которым у них нет прав.
    \item    Уьбедиться, что пользователи не могут обойти процесс авторизации.
\end{compactitem}

\textbf{Методы тестирования}


Существует множество методов, которые можно использовать для проверки авторизации и регистрации. Некоторые из наиболее распространенных методов включают в себя:

\begin{compactitem}

    \item    Тестирование «черного ящика»: при тестировании «черного ящика» тестер не имеет никаких знаний о внутренней реализации программного обеспечения. Тестер просто вводит данные в программу и наблюдает за выводом.
    \item    Тестирование белого ящика: при тестировании белого ящика тестер знает внутреннюю реализацию программного обеспечения. Эти знания можно использовать для выявления потенциальных уязвимостей безопасности.
    \item    Тестирование безопасности. Тестирование безопасности — это тип тестирования, направленный на выявление и устранение уязвимостей безопасности. Тестирование безопасности может использоваться для выявления уязвимостей в процессе авторизации и регистрации.
    
\end{compactitem}

\textbf{Инструменты тестирования}

Существует множество инструментов, которые можно использовать для автоматизации проверки авторизации и регистрации. Некоторые из самых популярных инструментов включают в себя:

\begin{compactitem}

    \item    OWASP Zed Attack Proxy (ZAP): ZAP — это бесплатный инструмент тестирования безопасности с открытым исходным кодом, который можно использовать для выявления и устранения уязвимостей безопасности.
    \item    Burp Suite: Burp Suite — это коммерческий инструмент для тестирования безопасности, который предлагает широкий спектр функций для выявления и устранения уязвимостей безопасности.
    \item    Nessus: Nessus — это коммерческий сканер уязвимостей, который можно использовать для выявления уязвимостей безопасности в различных системах.

\end{compactitem}

\textbf{Контрольный список тестирования}

Ниже приведен контрольный список пунктов, которые следует учитывать при тестировании авторизации и регистрации:

\begin{compactitem}

    \item    Уьбедиться, что пользователи могут получить доступ только к тем ресурсам и функциям, на доступ к которым они имеют право.
    \item    Уьбедиться, что пользователи не могут получить доступ к ресурсам и функциям, на доступ к которым у них нет прав.
    \item    Уьбедиться, что пользователи не могут обойти процесс авторизации.
    \item    Уьбедиться, что процесс авторизации безопасен и его нельзя легко обойти.
    \item    Уьбедиться, что процесс авторизации является масштабируемым и может обрабатывать большое количество пользователей.
    \item    Уьбедиться, что процесс авторизации надежен и не дает сбоев под нагрузкой.

\end{compactitem}

\textbf{Заключение}

Тестирование авторизации и регистрации является важной частью тестирования программного обеспечения. Тестируя процесс авторизации и регистрации, тестировщики могут помочь обеспечить безопасность и конфиденциальность данных пользователей.



\end{document}
